\section{Stage~L1 -- Coarse Architectural Boundary Selection}
\label{sec:rq11}

\subsection{Stage Objective}
As the first controlled-local stage of RQ1, this stage addresses the following question:

\begin{quote}
\textit{Which coarse architectural stage boundaries provide the most suitable two-part microservice decompositions of ResNet-18 under controlled local execution, and how do their latency and boundary-crossing costs compare to monolithic inference?}
\end{quote}

Stage~L1 is designed as a coarse-boundary screening stage rather than an exhaustive exploration of all possible split points. Its purpose is to identify architecturally meaningful two-part decompositions that can be compared against monolithic inference under controlled local execution, and to determine which coarse boundaries provide the strongest carry-forward basis for later refinement. In this stage, suitability is not treated as a question of raw latency alone. It is evaluated in relation to the joint trade-off between end-to-end latency, activation-transfer burden, boundary-crossing cost, and the structural balance of computation across the two resulting services.

The literature framing for coarse architectural boundary selection is discussed
in \cref{sec:rw-split-computing,sec:rw-granularity}.

\subsection{Experimental Setup}
The Stage~L1 benchmark used the fully controlled local CPU-only configuration in order to minimize variance and isolate the effect of the partition boundary itself. The source configuration file is \texttt{configs/rq1/1.1/rq1\_1\_fully\_controlled.yaml}; an exact byte-for-byte snapshot is preserved alongside the frozen results as \texttt{config.yaml} in the artifact directory. The evaluated model was ResNet-18 with pretrained weights, executed on CPU in FP32 precision with an input shape of $1 \times 3 \times 224 \times 224$. The benchmark compared one monolithic baseline against four two-part split configurations, each corresponding to a coarse architectural stage boundary in ResNet-18.

Activation-tensor sizes in this section are quoted in binary kibibytes (KiB, $1024$ bytes) for consistency with the gRPC message-size limit, which is also a binary quantity. Three further configuration files are present under \texttt{configs/rq1/1.1/} but do not contribute to the local-controlled Stage~L1 thesis-facing result reported here: \texttt{rq1\_1\_experimental.yaml} (a smoke test with minimal iterations, used only for pipeline validation), \texttt{rq1\_1\_threaded.yaml} (a four-core ancillary profile retained as a historical comparison point), and \texttt{rq1\_1\_azure\_fully\_controlled.yaml} (an AKS-backed coarse-split variant driven by \texttt{scripts/run\_rq11\_azure\_fully\_controlled.py}, retained for completeness but not part of the local single-host result reported in this section).

\begin{table}[htbp]
\centering
\caption{Stage~L1 benchmark configuration}
\label{tab:rq11-config}
\begin{tabular}{ll}
\toprule
\textbf{Setting} & \textbf{Value} \\
\midrule
Model & ResNet-18 (pretrained) \\
Device & CPU \\
Precision & FP32 \\
Input shape & $1 \times 3 \times 224 \times 224$ \\
Conditions & Monolithic + split after layer1, layer2, layer3, layer4 \\
Rounds & 5 \\
Warmup iterations & 50 \\
Measured iterations & 200 \\
Cooldown between conditions & 5 seconds \\
Random seed & 42 \\
Communication & gRPC over localhost (\texttt{127.0.0.1:50051}) \\
Maximum message size & 16 MiB ($16777216$ bytes) \\
Parity tolerance & $1.0 \times 10^{-5}$ \\
Parity inputs & 5 \\
Warmup window & 10 iterations \\
Warmup CV threshold & 0.02 \\
\bottomrule
\end{tabular}
\end{table}

The benchmark was configured to use a single-threaded PyTorch execution model pinned to one physical CPU core with SMT avoided. Thread counts for PyTorch, OpenMP, MKL, and OpenBLAS were all fixed to one. In addition, process priority was increased, the CPU governor was set to \texttt{performance}, and turbo boost was disabled in order to reduce frequency-related variance. These controls were applied symmetrically across all conditions so that the comparison remained fair and reproducible.

\subsection{Benchmark Design}
The experiment followed a repeated-measures design with five conditions: a monolithic baseline and four split configurations. The split points were selected after the major architectural stages of ResNet-18, namely after \texttt{layer1}, \texttt{layer2}, \texttt{layer3}, and \texttt{layer4}. These boundaries are meaningful because they preserve the stage structure of the network while allowing the evaluation of progressively later partition points.

The benchmark protocol consisted of 5 rounds. For each round and condition, 50 warmup iterations were executed before 200 measured iterations were recorded. A 5-second cooldown separated successive conditions. The ordering of conditions was randomized using deterministic per-round seeds derived from the configured seed and round index to reduce ordering bias. Functional equivalence between monolithic and split executions was enforced as a mandatory precondition through both local and gRPC parity checks using five inputs and an absolute tolerance of $1.0 \times 10^{-5}$.

A warmup calibration procedure was used to monitor whether latency had stabilized. This procedure used a trailing window of 10 iterations and a coefficient-of-variation threshold of 0.02. The implementation could extend warmup adaptively beyond the configured minimum if the trailing window remained unstable; in this frozen run, all condition-round pairs stabilized within the configured 50 iterations, so no extra warmup iterations were used.

\subsection{Partition Conditions}
The four split conditions corresponded to the following coarse stage boundaries in ResNet-18:

\begin{itemize}
    \item \textbf{split\_after\_layer1:} The first service executes the stem and \texttt{layer1}, and the second service executes \texttt{layer2} through the classifier head. This split transfers the largest intermediate tensor, of exact shape $1 \times 64 \times 56 \times 56$, corresponding to $802{,}816$ bytes ($784$ KiB) in FP32.
    \item \textbf{split\_after\_layer2:} The first service executes the stem and \texttt{layer1--layer2}, and the second service executes \texttt{layer3} through the classifier head. The intermediate tensor is of exact shape $1 \times 128 \times 28 \times 28$, $401{,}408$ bytes ($392$ KiB).
    \item \textbf{split\_after\_layer3:} The first service executes the stem and \texttt{layer1--layer3}, and the second service executes \texttt{layer4} plus pooling and classification. The intermediate tensor is of exact shape $1 \times 256 \times 14 \times 14$, $200{,}704$ bytes ($196$ KiB).
    \item \textbf{split\_after\_layer4:} The first service executes the stem and \texttt{layer1--layer4}, and the second service executes only average pooling and the final linear layer. The intermediate tensor is of exact shape $1 \times 512 \times 7 \times 7$, $100{,}352$ bytes ($98$ KiB).
\end{itemize}

These stage boundaries provide a coarse sweep across the network from earlier to later partition points. This design is consistent with partitioning research that treats split placement as a search-space and granularity decision shaped by computation, communication, and deployment constraints, rather than as an unconstrained exhaustive sweep over every internal operation \cite{dnnsplit, pareto_split, survey_partitioning}.

\paragraph{Carry-forward rule.} The decision of which coarse boundary to promote to the next stage is governed by a predeclared two-part rule applied to the per-condition mean end-to-end latencies and the per-service compute measurements (see \cref{sec:methods-carryforward} for the full statement). In short: (i) the \emph{near-best window} comprises every split whose mean latency lies within $5\%$ of the raw-fastest split; (ii) a split is treated as \emph{compute-degenerate} when the smaller of the two per-service mean compute times contributes less than $10\%$ of the total split compute, i.e.\ $\min(A,B)/(A+B) < 0.10$ for service-A compute $A$ and service-B compute $B$. The main carry-forward candidate is the non-degenerate split with the lowest mean latency in the near-best window; the next-best non-degenerate split is retained as the reference candidate. The frozen rule outputs for this run are recorded in \texttt{carry\_forward.json}.

\paragraph{Boundary-crossing interval.} The \emph{boundary-crossing interval} reported for each split condition is the client-observed interval that begins immediately before the intermediate-tensor serialisation and ends immediately after the response deserialisation. It is a composite of activation serialisation, gRPC transport, the remote service-B compute, and response deserialisation, and is implemented by the timer block in \texttt{src/client/split\_client.py} (see \cref{sec:methods-contract} for the full measurement contract). It is not a pure ``wire'' cost.

\subsection{Results}
\label{sec:rq11-results}
The thesis-facing Stage~L1 benchmark used the frozen artifact
\texttt{rq1\_1\_20260515\_111428} (source:
\texttt{results/frozen\_new/rq1\_1\_20260515\_111428/condition\_summaries.csv}).
Monolithic execution achieved the lowest mean end-to-end latency at
77.24~ms. The four split conditions yielded mean latencies of 80.88~ms
for \texttt{split\_after\_layer1}, 80.12~ms for
\texttt{split\_after\_layer2}, 79.08~ms for \texttt{split\_after\_layer3},
and 79.58~ms for \texttt{split\_after\_layer4}. This corresponds to
relative overheads of approximately 4.7\%, 3.7\%, 2.4\%, and 3.0\%,
respectively, compared to monolithic execution.

\begin{table}[htbp]
\centering
\caption{Stage~L1 summary of mean latency, dispersion, descriptive within-run per-iteration 95\% confidence interval on the mean, and overhead vs.\ monolithic. Source:
\texttt{condition\_summaries.csv} in
\texttt{rq1\_1\_20260515\_111428}. Confidence intervals are reported from the same artifact's \texttt{ci95\_lower\_ms} / \texttt{ci95\_upper\_ms} columns and summarise per-iteration variation within this run.}
\label{tab:rq11-results}
\begin{tabular}{lrrrrr}
\toprule
\textbf{Condition} & \textbf{Mean (ms)} & \textbf{Median (ms)} & \textbf{Std. Dev. (ms)} & \textbf{Within-run 95\% CI (ms)} & \textbf{Overhead vs. Monolithic} \\
\midrule
Monolithic & 77.24 & 75.29 & 3.59 & $[77.01,\ 77.46]$ & -- \\
Split after layer1 & 80.88 & 80.24 & 2.27 & $[80.74,\ 81.02]$ & 4.7\% \\
Split after layer2 & 80.12 & 79.60 & 2.40 & $[79.97,\ 80.27]$ & 3.7\% \\
Split after layer3 & 79.08 & 78.38 & 2.31 & $[78.94,\ 79.23]$ & 2.4\% \\
Split after layer4 & 79.58 & 78.85 & 2.43 & $[79.43,\ 79.73]$ & 3.0\% \\
\bottomrule
\end{tabular}
\end{table}

\begin{figure}[htbp]
\centering
\resizebox{0.92\textwidth}{!}{%
% Auto-generated from results/frozen_new/rq1_1_20260515_111428/condition_summaries.csv and carry_forward.json
% via scripts/render_rq11_figure.py. Do not edit by hand -- re-run the script.
\begin{tikzpicture}[
    font=\small,
    axis/.style={thick, draw=black!70, -{Latex[length=2mm]}},
    grid/.style={draw=black!12, thin},
    baselinebar/.style={draw=black!60, fill=black!10, rounded corners=2pt, thick},
    selectedmainbar/.style={draw=green!45!black, fill=green!22, rounded corners=2pt, thick},
    selectedrefbar/.style={draw=green!40!black, fill=green!10, rounded corners=2pt, thick},
    rejectedbar/.style={draw=red!70!black, fill=red!10, rounded corners=2pt, thick},
    degeneratebar/.style={draw=orange!75!black, fill=orange!15, rounded corners=2pt, thick},
    label/.style={font=\scriptsize, text=black!75, align=center},
    vallabel/.style={font=\scriptsize\bfseries, text=black!75},
    note/.style={font=\scriptsize, text=black!65, align=center},
    errbar/.style={draw=black!80, semithick, line cap=round}
]

% Frozen means +/- 1 SD (ms) used for this figure:
%   monolithic: 77.2375 +/- 3.5873
%   split_after_layer1: 80.8774 +/- 2.2730
%   split_after_layer2: 80.1176 +/- 2.4037
%   split_after_layer3: 79.0843 +/- 2.3075
%   split_after_layer4: 79.5831 +/- 2.4343
% Raw fastest split: split_after_layer3
% Selected main: split_after_layer3
% Selected reference: split_after_layer2
% Degenerate candidates: ['split_after_layer4']
% Rejected: ['split_after_layer1', 'split_after_layer4']

\draw[grid] (0,0.000) -- (9.300,0.000);
\node[anchor=east, font=\scriptsize, text=black!65] at (-0.15,0.000) {73};
\draw[grid] (0,0.636) -- (9.300,0.636);
\node[anchor=east, font=\scriptsize, text=black!65] at (-0.15,0.636) {74};
\draw[grid] (0,1.273) -- (9.300,1.273);
\node[anchor=east, font=\scriptsize, text=black!65] at (-0.15,1.273) {75};
\draw[grid] (0,1.909) -- (9.300,1.909);
\node[anchor=east, font=\scriptsize, text=black!65] at (-0.15,1.909) {76};
\draw[grid] (0,2.545) -- (9.300,2.545);
\node[anchor=east, font=\scriptsize, text=black!65] at (-0.15,2.545) {77};
\draw[grid] (0,3.182) -- (9.300,3.182);
\node[anchor=east, font=\scriptsize, text=black!65] at (-0.15,3.182) {78};
\draw[grid] (0,3.818) -- (9.300,3.818);
\node[anchor=east, font=\scriptsize, text=black!65] at (-0.15,3.818) {79};
\draw[grid] (0,4.455) -- (9.300,4.455);
\node[anchor=east, font=\scriptsize, text=black!65] at (-0.15,4.455) {80};
\draw[grid] (0,5.091) -- (9.300,5.091);
\node[anchor=east, font=\scriptsize, text=black!65] at (-0.15,5.091) {81};
\draw[grid] (0,5.727) -- (9.300,5.727);
\node[anchor=east, font=\scriptsize, text=black!65] at (-0.15,5.727) {82};
\draw[grid] (0,6.364) -- (9.300,6.364);
\node[anchor=east, font=\scriptsize, text=black!65] at (-0.15,6.364) {83};
\draw[grid] (0,7.000) -- (9.300,7.000);
\node[anchor=east, font=\scriptsize, text=black!65] at (-0.15,7.000) {84};

\draw[axis] (0,0) -- (9.600,0);
\draw[axis] (0,0) -- (0,7.400);

\node[rotate=90, font=\small, text=black!75] at (-0.85,3.500) {Mean end-to-end latency (ms)};
\node[font=\small, text=black!75] at (5.050,-1.15) {Stage~L1 condition};

\draw[baselinebar] (0.800,0) rectangle (1.800,2.697);
\draw[rejectedbar] (2.600,0) rectangle (3.600,5.013);
\draw[selectedrefbar] (4.400,0) rectangle (5.400,4.529);
\draw[selectedmainbar] (6.200,0) rectangle (7.200,3.872);
\draw[degeneratebar] (8.000,0) rectangle (9.000,4.189);

\draw[errbar] (1.300,0.414) -- (1.300,4.979);
\draw[errbar] (1.120,4.979) -- (1.480,4.979);
\draw[errbar] (1.120,0.414) -- (1.480,0.414);
\draw[errbar] (3.100,3.566) -- (3.100,6.459);
\draw[errbar] (2.920,6.459) -- (3.280,6.459);
\draw[errbar] (2.920,3.566) -- (3.280,3.566);
\draw[errbar] (4.900,3.000) -- (4.900,6.059);
\draw[errbar] (4.720,6.059) -- (5.080,6.059);
\draw[errbar] (4.720,3.000) -- (5.080,3.000);
\draw[errbar] (6.700,2.403) -- (6.700,5.340);
\draw[errbar] (6.520,5.340) -- (6.880,5.340);
\draw[errbar] (6.520,2.403) -- (6.880,2.403);
\draw[errbar] (8.500,2.640) -- (8.500,5.738);
\draw[errbar] (8.320,5.738) -- (8.680,5.738);
\draw[errbar] (8.320,2.640) -- (8.680,2.640);

\node[vallabel, above] at (1.300,4.979) {77.24};
\node[vallabel, above] at (3.100,6.459) {80.88};
\node[vallabel, above] at (4.900,6.059) {80.12};
\node[vallabel, above] at (6.700,5.340) {79.08};
\node[vallabel, above] at (8.500,5.738) {79.58};

\node[label] at (1.300,-0.45) {\texttt{mono}};
\node[label] at (3.100,-0.45) {\texttt{layer1}\\rejected};
\node[label] at (4.900,-0.45) {\texttt{layer2}\\reference};
\node[label] at (6.700,-0.45) {\texttt{layer3}\\selected};
\node[label] at (8.500,-0.45) {\texttt{layer4}\\degenerate};

\node[note, text width=9.2cm] at (5.050,-1.85) {
    \texttt{split\_after\_layer3} is the raw-fastest split. \texttt{split\_after\_layer4} is excluded by the compute-degeneracy rule (minor side $<$ 10\% of split compute). The carry-forward rule (5\% near-best window) selects \texttt{split\_after\_layer3} as the main candidate and \texttt{split\_after\_layer2} as the reference candidate.
};

\end{tikzpicture}

}
\caption{Stage~L1 coarse split latency results under controlled local execution. Bar heights are the mean end-to-end latency per condition for the frozen run \texttt{rq1\_1\_20260515\_111428}; bar colours encode the role assigned by the carry-forward rule (baseline, selected\_main, selected\_reference, rejected, degenerate). \texttt{split\_after\_layer3} is the raw-fastest split in this run and is also non-degenerate; \texttt{split\_after\_layer4} is excluded by the compute-degeneracy rule despite its small intermediate activation. The carry-forward decision (see \texttt{carry\_forward.json} for the same run) selects \texttt{split\_after\_layer3} as the main candidate and \texttt{split\_after\_layer2} as the reference candidate. Error bars show $\pm 1$ standard deviation of per-iteration latency (the \texttt{std\_ms} column of \cref{tab:rq11-results}); they are wider than the gaps between the non-degenerate split means, which is consistent with the scope caveat that the ranking is a within-run ordering rather than a host-independent separation (the much tighter within-run 95\% confidence interval on the mean is given in \cref{tab:rq11-results}). The figure is generated deterministically from the frozen \texttt{condition\_summaries.csv} and \texttt{carry\_forward.json} by \texttt{scripts/render\_rq11\_figure.py}.}
\label{fig:rq11_coarse_split_results}
\end{figure}

\Cref{fig:rq11_coarse_split_results} visualises the latency pattern reported in
\cref{tab:rq11-results}. In this frozen run the raw-fastest split is the
mid-to-late boundary after \texttt{layer3}, not the latest boundary, and the
carry-forward decision still depends on both latency and structural compute
balance rather than latency alone.

All split configurations therefore introduced positive latency overhead relative to monolithic inference. However, the magnitude and character of this overhead varied substantially with split location. The earliest split, after \texttt{layer1}, performed worst because it crossed the largest intermediate activation and exhibited the highest boundary-crossing cost. The latest split, after \texttt{layer4}, minimized activation-transfer size but left only a minimal residual computation segment on the second service (well under 1~ms of remote compute according to the per-service compute means in \texttt{condition\_summaries.csv}), making it structurally compute-degenerate despite its low transfer burden. The mid-to-late splits, after \texttt{layer2} and \texttt{layer3}, formed the strongest non-degenerate candidates in the coarse sweep, with \texttt{split\_after\_layer3} also yielding the lowest raw split mean in this frozen run.

The boundary-crossing measurements further support this interpretation. The split after \texttt{layer1} incurred the largest boundary-crossing interval at 53.81~ms, followed by \texttt{layer2} at 40.85~ms and \texttt{layer3} at 26.06~ms, while \texttt{layer4} was much smaller at 2.22~ms. This descending pattern is consistent with the joint reduction in activation-transfer burden and downstream service-B compute as the boundary moves later, from $784$~KiB at \texttt{layer1} to $98$~KiB at \texttt{layer4}. At the same time, the late \texttt{layer4} split shows that minimizing transfer burden alone does not determine overall suitability as a microservice decomposition, because the downstream service becomes too small to represent a meaningful balanced distribution.

Cross-round consistency was strong across all configurations, indicating stable execution behavior. Per \texttt{round\_consistency.json} for the same artifact, the monolithic condition had a round coefficient of variation (CV) of 0.002, while \texttt{split\_after\_layer1}, \texttt{split\_after\_layer2}, \texttt{split\_after\_layer3}, and \texttt{split\_after\_layer4} had round CVs of 0.005, 0.011, 0.003, and 0.005, respectively.

\paragraph{Scope caveat.} The numbers reported in this section are derived from a single frozen artifact (\texttt{rq1\_1\_20260515\_111428}) generated by a single benchmark run with the fixed seed $42$. The five rounds within that run randomise condition order but do not vary the input seed, the host, or the build, and cross-host or cross-seed replication is not in scope for Stage~L1. The mean-latency differences between the non-degenerate split conditions are small relative to their per-iteration standard deviation. Although the descriptive within-run per-iteration 95\% confidence intervals in \cref{tab:rq11-results} do not overlap between \texttt{split\_after\_layer3} and \texttt{split\_after\_layer2}, those intervals should not be read as host-independent uncertainty estimates. The absolute ranking within the non-degenerate set should therefore be read as ``the strongest carry-forward candidate observed in this frozen run'' rather than as a general ordering. The carry-forward decision is stated in terms of the predeclared rule (\cref{sec:methods-carryforward}) applied to this artifact, and the ranking is carried forward for refinement in Stage~L2 rather than treated as a final answer.

The interpretation of these results and their contribution to RQ1 are presented in
\cref{sec:analysis-rq11}.
